\section{Random energy: a joint law without a low-order graph}

Derrida's random-energy model assigns an independent energy to every complete configuration
\cite{derrida1980}.  For a $q$-state sequence of length $L$, take
\begin{equation}
E(x)\overset{\mathrm{iid}}{\sim}\mathcal N(0,L\sigma^2),
\qquad x\in[q]^L,
\qquad
P_\beta(x)=\frac{e^{-\beta E(x)}}{Z_L(\beta)}.
\label{eq:book-rem-law}
\end{equation}
The model has one exact joint Gibbs law.  Its low-temperature freezing transition follows from
competition between the $q^L$ configurations and the extreme tail of their energies.  With the
normalization in Eq.~\ref{eq:book-rem-law}, the critical inverse temperature is
$\beta_c=\sqrt{2\log q}/\sigma$.  Positional locality plays no role in this result.

The operator framework reveals where the random energy resides in interaction order.  Use the
uniform product reference and let $\Pi_A$ be the functional-ANOVA projector for a subset
$A\subseteq[L]$.

\begin{proposition}[ANOVA order profile of a random energy table]
\label{prop:book-rem-anova}
For the ensemble in Eq.~\ref{eq:book-rem-law},
\begin{equation}
\mathbb E_E\|\Pi_AE\|_{L^2}^2
=\frac{L\sigma^2}{q^L}(q-1)^{|A|}.
\label{eq:book-rem-subset-mass}
\end{equation}
Consequently the expected mass at interaction order $k$ is
\begin{equation}
W_k
=\frac{L\sigma^2}{q^L}\binom{L}{k}(q-1)^k,
\qquad
\sum_{k=0}^{L}W_k=L\sigma^2.
\label{eq:book-rem-order-profile}
\end{equation}
After normalization by total mass, the order profile is the binomial law
$\operatorname{Binomial}(L,(q-1)/q)$.
\end{proposition}

\begin{proof}
The product ANOVA subspace indexed by $A$ has dimension $(q-1)^{|A|}$.  The random vector of
$q^L$ energy values is isotropic with coordinate variance $L\sigma^2$.  The expected Euclidean
squared norm of its projection onto a rank-$r$ subspace is $L\sigma^2r$.  The uniform $L^2$ norm
divides by $q^L$, giving Eq.~\ref{eq:book-rem-subset-mass}.  Summing over the
$\binom{L}{k}$ subsets of size $k$ gives Eq.~\ref{eq:book-rem-order-profile}; the binomial theorem
gives the final identity.
\end{proof}

For protein alphabets, $(q-1)/q$ is close to one, so the ensemble mass is concentrated at very
high order.  This is an exact statement about a random function on sequence space, independent of
whether its Gibbs measure is in the frozen phase.

The same calculation sharpens the static--dynamic distinction.  Fix a pair $(i,j)$ and double
center the $q\times q$ energy table obtained at each background
$C=x_{[L]\setminus\{i,j\}}$.  The expected
total conditional pair-response mass is
\begin{equation}
\mathbb E_E\mathbb E_C\|G_{i\leftarrow j}(C)\|^2
=\beta^2L\sigma^2\frac{(q-1)^2}{q^2},
\label{eq:book-rem-local-response}
\end{equation}
whereas the expected squared norm of the background mean is
\begin{equation}
\mathbb E_E\|\overline G_{i\leftarrow j}\|^2
=\beta^2L\sigma^2\frac{(q-1)^2}{q^L}.
\label{eq:book-rem-static-response}
\end{equation}
Thus the ratio of expected static mass to expected total response mass is $q^{-(L-2)}$.  A typical background displays a substantial
finite pair response, while averaging across backgrounds removes almost all of it.  The apparent
edge is a section through a high-order landscape.

The random-energy model is therefore a useful null for protein sequence models.  It preserves a
globally normalized energy landscape and destroys reusable positional structure.  A learned model
has gone beyond this null when low-order operator mass, cross-background singular directions, or
family-transferable static blocks exceed matched random-function and sequence-shuffling controls.
Freezing, low-order interaction, and biological constraint remain separate empirical coordinates.
